\problema{Alien Dictionary}{269}{src/01-grafos/269-alien-dictionary.cpp}{H}

Nos dan una lista de palabras que ya están ordenadas según las reglas de un
alfabeto alienígena que no conocemos. A partir de esas palabras hay que
deducir un orden válido de las letras del alfabeto. Si la lista es
inconsistente (no existe ningún orden posible), devolvemos \texttt{""}.

\paragraph{La idea, con un ejemplo de la vida real.} Es como jugar a adivinar
el orden del abecedario de otro planeta usando su diccionario. Sabemos que las
palabras del diccionario ya están ordenadas alfabéticamente, aunque no sabemos
en qué orden van sus letras. ¿Cómo se compara dos palabras en un diccionario?
Se leen letra por letra hasta encontrar la primera que sea distinta; esa
primera diferencia decide cuál palabra va antes.

Entonces, cada par de palabras \emph{vecinas} (una y la que le sigue) nos regala
como mucho una pista: la primera letra en la que se diferencian nos dice ``esta
letra va antes que esa otra''. Las demás letras no dicen nada sobre el orden.
Con todas esas pistas dibujamos flechas letra $\to$ letra (un \emph{grafo
dirigido}: puntos que son letras y flechas que dicen quién va antes). Después
solo hay que acomodar las letras en una fila que respete todas las flechas; a
esa fila se le llama \emph{orden topológico}, y esa es la respuesta.

\begin{center}
\ttfamily
\begin{tabular}{lll}
wrt & wrf & \small (diff. índice 2: t $\to$ f)\\
wrf & er & \small (diff. índice 0: w $\to$ e)\\
er & ett & \small (diff. índice 1: r $\to$ t)\\
ett & rftt & \small (diff. índice 0: e $\to$ r)\\
\end{tabular}

\medskip
Grafo resultante:\quad w $\to$ e $\to$ r $\to$ t $\to$ f
\end{center}

\noindent En este ejemplo las aristas forman una sola cadena, así que el
orden topológico es único: \texttt{"wertf"}.

\paragraph{Cómo lo resuelve el código, paso a paso.} Primero juntamos las
pistas, y luego ordenamos las letras. El \emph{grado de entrada} de una letra es
cuántas flechas le llegan, es decir, cuántas letras deben ir antes que ella; una
letra con grado de entrada $0$ no tiene ninguna letra obligada delante.
\begin{enumerate}
  \item Comparamos cada palabra con la que le sigue y buscamos la primera letra
        en la que se diferencian: esa es la única pista (la única flecha) que
        aporta ese par.
  \item Con todas las flechas juntas aplicamos el \emph{algoritmo de Kahn}:
        metemos en una \emph{cola} (la fila del súper: primero en entrar,
        primero en salir) las letras con grado de entrada $0$.
  \item Sacamos una letra, la ponemos en la respuesta y a las letras que iban
        después de ella les quitamos esa flecha; la que se quede con grado de
        entrada $0$ entra a la cola.
  \item Repetimos hasta vaciar la cola.
\end{enumerate}
Si al final acomodamos todas las letras, esa fila es el orden del alfabeto. Si
sobró alguna letra (quedó atrapada en un círculo vicioso, un \emph{ciclo}, del
tipo ``a antes que b, b antes que a''), no existe orden posible y devolvemos
\texttt{""}.

Ojo con una trampa fácil de olvidar: si dos palabras vecinas empiezan igual pero
la primera es \emph{más larga} que la segunda (por ejemplo
\texttt{["abc", "ab"]}), eso rompe el orden del diccionario, porque una palabra
corta siempre debe ir antes que la palabra larga que empieza igual. En ese caso
devolvemos \texttt{""} de inmediato, sin siquiera ponernos a ordenar letras.

\lstinputlisting{src/01-grafos/269-alien-dictionary.cpp}

\complejidad{$O(C)$}{$O(1)$}

\paragraph{¿Qué significa esa complejidad?} Aquí $C$ es la suma de las
longitudes de todas las palabras, es decir, cuántas letras hay que leer en total
para comparar todas las palabras vecinas. Ese conteo es lo que manda en el
tiempo, porque acomodar las letras es rapidísimo: el alfabeto tiene a lo más
$26$ letras. Es lo más rápido que se puede, ya que hay que leer las palabras al
menos una vez para sacar las pistas.

\paragraph{Consejos para la entrevista (Amazon).} Es pariente cercano de Course
Schedule (207/210): usa el mismo algoritmo de Kahn, pero aquí lo difícil es
armar bien el dibujo de flechas a partir de las palabras. Los dos errores más
comunes: (1) olvidar la trampa de la palabra corta que va después de la larga
que empieza igual, y (2) tomar más de una letra distinta por cada par de
palabras (solo importa la \emph{primera} diferencia; lo demás no dice nada).
Recuerda también que casi siempre hay varios órdenes válidos (una letra que
nadie obliga a nada puede ir en varios lugares), así que conviene aclarar que
basta con devolver \emph{un} orden que funcione, no ``el'' orden.
